\section{权限与安全：在效率和安全之间找平衡}

\subsection{五种权限模式}

Claude Code 提供五种权限模式，用 \texttt{Shift+Tab} 在对话中循环切换，从全自动到完全手动确认，覆盖不同的安全需求场景。

\subsection{权限规则配置}

在 \texttt{.claude/settings.json} 里精细控制每种工具的权限。

\begin{importantbox}{权限优先级：deny > ask > allow}
可以大胆地 allow 常用命令，同时用 deny 保护敏感文件。比如 \texttt{.env} 文件无论如何都不应该被编辑。
\end{importantbox}

设置文件也有多级层次：团队共享规则放 \texttt{.claude/settings.json}，个人偏好放 \texttt{\textasciitilde/.claude/settings.json}，敏感配置放 \texttt{.claude/settings.local.json}（加到 \texttt{.gitignore}）。

\subsection{Hooks：让规则变成铁律}

\begin{warningbox}{CLAUDE.md 是"建议"，Hooks 是"铁律"}
\texttt{CLAUDE.md} 里的指令 Claude 大部分时候会遵守，但偶尔会忘。Hooks 是在特定生命周期事件上自动触发的 shell 命令，无论如何都会执行。配置在 \texttt{settings.json} 里。
\end{warningbox}

Hooks 的关键事件类型包括编辑文件前/后、执行命令前/后、上下文压缩后等。Hook 命令返回 \texttt{exit code 0} 表示允许，\texttt{exit code 2} 表示阻止，因此可以写任意复杂的判断逻辑。

实用场景举例：

\begin{itemize}[leftmargin=1.5em]
  \item \textbf{自动格式化}：每次编辑后运行 Prettier
  \item \textbf{保护关键文件}：阻止修改生产配置
  \item \textbf{压缩后重注入}：对话被压缩后自动重新注入关键指令，解决"长对话压缩丢信息"的痛点
\end{itemize}

Hook 有四种类型：\texttt{command}（shell 命令）、\texttt{http}（POST 请求）、\texttt{prompt}（单次 LLM 判断）、\texttt{agent}（启动子 Agent 验证）。

\subsection{本章小结}

权限管理的核心是分层：deny 保护底线，allow 提升效率。Hooks 把 CLAUDE.md 中的"建议"变成可执行的"铁律"，是保障代码安全的最后一道防线。

%% ====================================================================
